\section{导数是最佳线性映射}
\label{sec:10-Matrix-Calculus/01-Derivative-Conventions/01}

你给一个自定义算子写 \texttt{backward}，框架要求返回的东西与输入同形状，是一个 \(B\times d\) 的张量。同一周读到的控制论文里，``导数''是一个 \(6\times 7\) 的矩阵。再翻回一元微积分的讲义，导数是一个数。三处用的是同一个词，写下来的对象从标量长到向量再长到矩阵，彼此连能不能相加都要先看形状。

一种应付办法是把几套公式分开背，各管一摊：标量对标量、标量对向量、向量对向量。教科书里的例子这样确实能对付过去。可只要碰上没背过的组合——输入是矩阵、输出是向量，或者两端都是矩阵——手上就没有模板了，而深度学习里这种组合是常态：一个全连接层的权重是矩阵，损失是标量，中间还夹着若干张量。公式表越背越长，判断力却没有增长。

另一条路是承认这三样东西记录的是同一个对象。给定一点 \(x\)，函数在它附近能被一个线性映射近似到一阶精度；那个线性映射就是导数。它把``输入的一个微小扰动''送到``输出相应的微小变化''。输入扰动住在什么空间、输出变化住在什么空间，决定了我们最后把这个映射写成几行几列的数组；映射本身不随记号改变。

\begin{center}
\begin{tikzpicture}[x=1cm,y=1cm,font=\small,>=stealth]
  \node[draw=blue!55!black,thick,rounded corners=3pt,inner sep=6pt,align=center,fill=blue!5]
    at (4.6,4.35) {同一个对象：线性映射 \(Df(x):\ h\ \longmapsto\ Df(x)[h]\)};
  \draw[->,black!55] (3.1,3.98) -- (1.3,2.80);
  \draw[->,black!55] (4.6,3.98) -- (4.6,2.74);
  \draw[->,black!55] (6.1,3.98) -- (7.9,3.22);
  % 一维：一个格子
  \draw[black!70,fill=blue!8] (1.09,2.19) rectangle (1.51,2.61);
  % 标量输出：一行 n 个格子
  \foreach \i in {0,1,2,3}
    \draw[black!70,fill=blue!8] (3.76+0.42*\i,2.19) rectangle (4.18+0.42*\i,2.61);
  % 一般情形：m 行 n 列
  \foreach \i in {0,1,2,3}
    \foreach \j in {0,1,2}
      \draw[black!70,fill=blue!8] (7.16+0.42*\i,1.77+0.42*\j) rectangle (7.58+0.42*\i,2.19+0.42*\j);
  \node[font=\footnotesize] at (1.30,1.32) {\(f:\R\to\R\)};
  \node[font=\footnotesize] at (1.30,0.82) {\(f'(x)\)：\(1\) 个数};
  \node[font=\footnotesize] at (4.60,1.32) {\(f:\R^n\to\R\)};
  \node[font=\footnotesize] at (4.60,0.82) {\(\nabla f(x)\)：\(n\) 个数};
  \node[font=\footnotesize] at (8.04,1.32) {\(f:\R^n\to\R^m\)};
  \node[font=\footnotesize] at (8.04,0.82) {\(J_f(x)\)：\(m\times n\) 个数};
\end{tikzpicture}

\emph{下面三排格子是同一个映射的三份账本；格子怎么排随定义域变化，格子上方那个箭头不变。}
\end{center}

图上方那个框是本单元的主线。导数首先是一个线性映射，矩阵、向量、单个数都只是它在选定基下的记账方式。这条判断在线性代数里已经出现过一次，见\hyperref[sec:03-Linear-Algebra/02-Vector-Spaces/04]{``线性变换与换基：对象不变，坐标可以改变''}；矩阵微积分把它原样搬到求导上。先认下这一条，后面所有关于布局、微分法、链式法则顺序的讨论才有一个共同的落脚点，否则每一条规则都要当成新知识去记。

\subsection{一维的导数早就是一个线性映射}

回头端详最熟悉的情形。设 \(f:\R\to\R\) 在 \(x\) 处可导，把极限定义重排一下：

\[
f(x+h)=f(x)+f'(x)\,h+r(h),\qquad \frac{\abs{r(h)}}{\abs{h}}\to 0\ (h\to 0).
\]

三样东西被分开了。\(f(x)\) 是当前位置。\(f'(x)h\) 是一阶修正，它对 \(h\) 是线性的——这才是导数出场的地方。\(r(h)\) 是余项，它不只要趋于零，还要比 \(h\) 趋于零得更快，这个``更快''是全部要求所在。

\(f'(x)h\) 表面上是一次乘法。换个角度读，它是映射 \(h\mapsto f'(x)h\) 作用在 \(h\) 上的结果，只不过一维空间上的线性映射恰好都能写成``乘一个数''，于是那个数被误当成了导数本身。

把维度提上去，结构照搬。对 \(f:\R^n\to\R^m\)，扰动 \(h\) 是 \(n\) 维向量，输出变化是 \(m\) 维向量，中间需要一个从 \(\R^n\) 到 \(\R^m\) 的线性映射。若存在这样的 \(Df(x)\)，使

\[
f(x+h)=f(x)+Df(x)[h]+r(h),
\qquad
\frac{\norm{r(h)}}{\norm{h}}\to 0\ (h\to 0),
\]

就称 \(f\) 在 \(x\) 处 Fréchet 可微，\(Df(x)\) 是它的导数。

方括号是有意写成这样的。\(Df(x)[h]\) 提醒你左边是一个作用在方向 \(h\) 上的映射，而不是一个被动装着偏导数的容器。选定坐标基之后，我们才把它落地成矩阵。这个区分不是咬文嚼字：后面会看到，一个函数可以在所有方向上都有方向导数，却根本不存在一个统一的线性映射把它们全部近似出来。先入为主地把导数等同于 Jacobian，跳过的恰好是``那个映射是否存在''这一步。

\subsection{Jacobian 是这个映射在标准基下的账本}

在标准基下，\(Df(x)\) 对应矩阵 \(J_f(x)\in\R^{m\times n}\)，使得 \(f(x+h)\approx f(x)+J_f(x)h\)。行数是输出维度，列数是输入维度。形状由这两个数定死，与你打算用什么求导公式无关。输出变成三维、输入保持二维，Jacobian 就必须是 \(3\times 2\)——这个判断在动手算第一个偏导数之前就能作出。

看一个具体的：

\[
f(x_1,x_2)=\begin{bmatrix} x_1x_2 \\ \sin x_1 \end{bmatrix}.
\]

令 \(h=(h_1,h_2)^{\trans}\)。第一个分量展开 \((x_1+h_1)(x_2+h_2)-x_1x_2=x_2h_1+x_1h_2+h_1h_2\)，末项是二阶量，扔进 \(r(h)\)。第二个分量用 \(\sin\) 的一阶展开得 \(\cos(x_1)h_1\)。于是

\[
Df(x)[h]=\begin{bmatrix} x_2h_1+x_1h_2 \\ \cos(x_1)h_1 \end{bmatrix},
\qquad
J_f(x)=\begin{bmatrix} x_2 & x_1 \\ \cos x_1 & 0 \end{bmatrix}.
\]

两行对应两个输出，两列对应两个输入。第 \(j\) 列回答的是：单独微调第 \(j\) 个输入，各个输出分量会挪动多少。

这个结构在工程里随处可见。调淋浴的冷热水旋钮，两个旋钮位置记作 \(q\)，你关心的输出 \(p(q)\) 是水温和总流量，一次微小调节满足 \(\Delta p\approx J_p(q)\Delta q\)。想在升温的同时保住流量，要解的是一个反问题：找 \(\Delta q\) 使 \(\Delta p\) 落在指定方向上。热水快断的时候，两个旋钮对水温都几乎不起作用，\(J_p(q)\) 接近秩亏，那个方向怎么拧也够不着。机器人手臂的末端控制是同一件事换了变量名：关节角代替旋钮，末端位姿代替水温流量，\(J\) 接近秩亏对应运动学奇异位形。工业机械臂在奇异位形附近降速，原因是此时 Jacobian 的伪逆会把一个小小的末端指令放大成巨大的关节速度，见\hyperref[sec:03-Linear-Algebra/07-SVD-and-Data-Applications/02]{``伪逆与病态问题''}。

共同的骨架是：一个多入多出的非线性映射，在局部被它的 Jacobian 线性近似；可微性保证这个矩阵存在，它的秩决定了哪些输出方向在当前点够得着。Jacobian 与 Fréchet 导数并非两件事，前者是后者穿上坐标外衣的样子。

\subsection{梯度是内积挑出来的那个向量}

输出是一维时，\(f:\R^n\to\R\)，\(Df(x)\) 把 \(n\) 维向量送成一个标量，是一个线性泛函。有限维欧氏空间里每个线性泛函都能唯一写成与某个向量作内积（Riesz 表示），于是

\[
Df(x)[h]=\nabla f(x)^{\trans}h,
\]

这个被挑出来的向量就叫梯度。它与 \(x\) 同形状：\(x\) 是 \(n\) 维列向量，\(\nabla f(x)\) 也是 \(n\) 维列向量。注意这里已经悄悄用了一次约定——线性泛函本身更像一个行向量，把它写成列向量是内积替我们做的翻译。下一节要讨论的布局分歧，根子就在这次翻译上。

走一个最常用的例子。取 \(f(x)=\tfrac12\norm{Ax-b}_2^2\)，展开到一阶：

\[
\begin{aligned}
f(x+h)&=\tfrac12\norm{Ax-b}_2^2+(Ax-b)^{\trans}Ah+\tfrac12\norm{Ah}_2^2\\
&=f(x)+\bigl[A^{\trans}(Ax-b)\bigr]^{\trans}h+r(h).
\end{aligned}
\]

第三项只含 \(h\) 的二次量，除以 \(\norm{h}\) 后随 \(h\to 0\) 消失，可以放进余项。中间那项对 \(h\) 线性，整理成``某个向量与 \(h\) 的内积''之后，那个向量就是梯度：\(\nabla f(x)=A^{\trans}(Ax-b)\)。整个过程没有出现任何一个偏导数符号。

梯度指向最速上升方向，这也不是外加的性质。由 Cauchy--Schwarz，\(\nabla f(x)^{\trans}h\le\norm{\nabla f(x)}\norm{h}\)，等号在 \(h\) 与 \(\nabla f(x)\) 同向时取到，所以在所有单位方向里它使一阶增量最大。线性泛函在单位球面上的极值方向，就是它的 Riesz 表示向量的方向。

有一个容易被跳过的推论。梯度这个向量取决于你选了哪个内积。设正定矩阵 \(M\) 定义内积 \(\ip{u}{v}_M=u^{\trans}Mv\)，对应的梯度 \(g_M\) 须满足 \(g_M^{\trans}Mh=\nabla f(x)^{\trans}h\) 对所有 \(h\) 成立，也就是 \(Mg_M=\nabla f(x)\)，于是 \(g_M=M^{-1}\nabla f(x)\)。导数映射一动没动，变的是拿什么几何把它表示成向量。自然梯度、预条件、Riemann 流形上的最速下降，共同的出处都是这一行等式；在参数空间弯曲的地方，Fisher 信息矩阵扮演的就是 \(M\)。优化里那些形如 \(M^{-1}\nabla f\) 的更新方向，与其说是对梯度做了修正，不如说是换了一把尺子重新读同一个导数，见\hyperref[sec:09-Convex-Optimization/06-Optimization-Algorithms/01]{``梯度下降真正难在步长''}。

\subsection{输入是矩阵时故事不变}

输入可以是矩阵。设 \(f:\R^{m\times n}\to\R\)，扰动 \(H\) 是一个同形状的矩阵，\(Df(X)[H]\) 仍是对 \(H\) 线性的标量。矩阵空间上的自然内积是 Frobenius 内积 \(\ip{A}{B}_F=\trace(A^{\trans}B)\)，用它把导数表示出来：

\[
Df(X)[H]=\ip{\nabla_X f}{H}_F=\trace\bigl((\nabla_X f)^{\trans}H\bigr),
\]

\(\nabla_X f\) 与 \(X\) 同形状。你不需要先把 \(X\) 拉平成长向量、求完导再折回矩阵，Frobenius 内积已经提供了现成的坐标表示。

最简单的例子：\(f(X)=\tfrac12\norm{X}_F^2\) 展开为 \(\tfrac12\norm{X}_F^2+\ip{X}{H}_F+\tfrac12\norm{H}_F^2\)，线性部分是 \(\ip{X}{H}_F\)，故 \(\nabla_X f=X\)。行列式、逆、迹的各种复合都走同一条路：展开到一阶，把含 \(H\) 的线性部分整理成 \(\trace(A^{\trans}H)\)，括号里的 \(A\) 就是梯度。这套操作在\hyperref[sec:10-Matrix-Calculus/02-Differentials/01]{``用微分与 Frobenius 内积读取梯度''}里会被立成一个可以反复执行的程序。

\subsection{链式法则只是把两个近似接起来}

若 \(g:\R^n\to\R^m\) 在 \(x\) 处可微、\(f:\R^m\to\R^p\) 在 \(g(x)\) 处可微，则复合 \(h=f\circ g\) 的导数是两个线性映射的复合：

\[
Dh(x)=Df(g(x))\circ Dg(x),
\qquad
J_h(x)=J_f(g(x))\,J_g(x).
\]

一阶近似再接一阶近似还是一阶近似，坐标写法就是 Jacobian 相乘。与一维的 \(h'(x)=f'(g(x))g'(x)\) 相比，唯一的变化是两个数的乘法换成了两个矩阵的乘法，而矩阵乘法不可交换，顺序被映射的作用次序钉死：先 \(Dg\)，后 \(Df\)。

这个顺序里藏着深度学习最要紧的一笔账。三个以上的 Jacobian 连乘时，结合律允许你从左往右乘，也允许从右往左乘，结果相同而算力开销可以差出一个维度。这条账在\hyperref[sec:10-Matrix-Calculus/02-Differentials/02]{``链式法则就是线性映射复合''}里单独算。

\subsection{方向导数拼不出全导数}

沿方向 \(h\) 的方向导数是 \(D_hf(x)=\lim_{t\to 0}\bigl(f(x+th)-f(x)\bigr)/t\)。若 Fréchet 导数存在，把 \(th\) 代入定义、除以 \(t\) 再取极限，用 \(\norm{r(th)}/t=\bigl(\norm{r(th)}/\norm{th}\bigr)\norm{h}\to 0\) 即得 \(D_hf(x)=Df(x)[h]\)，每个方向都有方向导数。

反过来不成立。所有方向的方向导数都存在，并不能保证有一个统一的线性映射同时近似它们。取

\[
f(x,y)=\frac{x^2y}{x^2+y^2},\qquad f(0,0)=0.
\]

原点处它连续（因为 \(\abs{x^2y/(x^2+y^2)}\le\abs{y}\)），沿任意方向 \(h=(a,b)\) 的方向导数都存在，等于 \(a^2b/(a^2+b^2)\)。检查线性性：\(D_{(1,1)}f(0)=\tfrac12\)，而 \(D_{(1,0)}f(0)+D_{(0,1)}f(0)=0\)。\(h\mapsto D_hf(0)\) 不是线性映射，可 Fréchet 可微会强迫它线性，所以全导数在原点不存在。更眼熟的例子是 \(f(x)=\norm{x}_2\) 在原点：沿每个方向的单侧斜率都等于 \(\norm{h}\)，而 \(\norm{h}\) 关于 \(h\) 不线性。

两个例子的教训一致。方向导数只检查函数沿一条条直线的行为，全导数要求在整个邻域内、对所有方向有一个联合近似，前者是后者的必要条件，离充分还很远。偏导数更窄，它只走坐标轴那几个方向，再多的坐标轴也铺不满一个邻域。

条件对一遍。若 \(f\) 在 \(x\) 的某个邻域内所有偏导数存在，且在 \(x\) 处连续，则 \(f\) 在 \(x\) 处 Fréchet 可微——这是从偏导升格到全导最常用的充分条件，连续偏导给出的局部控制足以挡住上面那种病态。反向不真：Fréchet 可微既不要求偏导连续，也不要求它们在 \(x\) 以外的地方存在。可微是一点上的性质，不对邻域作任何承诺。日常写代码时这些反例很少被撞上，因为常用的层由光滑运算搭成；真正需要停一下的是 \(\relu\)、\(\max\)、排序、取整这类地方，那里的``梯度''是一个约定而非导数。

回到开头那三个形状不同的对象。它们的差别现在有了统一的解释：定义域和陪域一变，同一个线性映射的账本就换一种排法。不过账本要排成什么样，还剩一个自由度没有定——梯度写成列还是写成行、Jacobian 是 \(m\times n\) 还是它的转置，数学上都说得通。下一节就来处理这个自由度，以及不把它挑明会付出什么代价。
